\lecture{5}{8. September 2025}{Rigid Body Equilibrium, Pt. 1}

\begin{problemwithimage}{f4_15}{F4-15}
  Determine the maximum force $P$ that can be applied without causing the two \qty{50}{\kg} crates to move.
  The coefficient of static friction between each crate and the ground is $\mu_s = \num{0,25}$.
\end{problemwithimage}

\begin{derivation}
  The maximum friction force that can be sustained by crate $A$ is:
  \begin{equation*}
    F_A = \mu_s W
  .\end{equation*}
  Crate $B$ is pulled upwards by the applied force, hence the friction limit for this is:
  \begin{equation*}
    F_B = \mu_s \cdot \left( W - P \cdot \sin \ang{30} \right)
  .\end{equation*}
  Hence, the horizontal force balance becomes:
  \begin{align*}
    P \cos \ang{30} &= F_A + F_B \\
    P \cos \ang{30} &= \mu_s W + \mu_s W - \mu_s P \cdot \sin \ang{30} \\
    P \left( \cos \ang{30} + \mu_s \sin \ang{30}  \right) &= 2 \mu_s W \\
    P &= \frac{2\mu_s W}{\cos \ang{30} + \mu_s \sin \ang{30} } \\
    &= \frac{2 \cdot \num{0.25} \cdot \qty{50}{\kg} \cdot \qty{9.81}{\m\per\second\squared} }{\cos \ang{30} + \num{0.25} \cdot \sin \ang{30} } \\
    &= \qty{247,42}{\N}
  .\end{align*}
\end{derivation}

\finalresult{P = \qty{247,42}{\N}}

\begin{problemwithimage}{f4_55}{4-55}
  The block brake consists of a pin-connected lever and friction block at $B$.
  The coefficient of static friction between the wheel and the lever is $\mu_s = \num{0,3}$, and a torque of \qty{5}{\N.\m} is applied to the wheel.
  Determine if the brake can hold the wheel stationary when the force applied to the lever is (a) $P = \qty{30}{\N}$, (b) $P = \qty{70}{\N}$.
\end{problemwithimage}

\begin{derivation}
  The tangential friction force needed to keep the wheel in place is:
  \begin{equation*}
    M = F_{\mathrm{req}} \cdot R \implies F_{\mathrm{req}} = \frac{M}{R} = \frac{\qty{5}{\N.\m} }{\qty{150}{\mm}} = \qty{33,33}{\N}
  .\end{equation*}
  This means the normal force $N$ at the block needs to be:
  \begin{equation*}
    F_{\mathrm{req}} = N \cdot \mu_s \implies N = \frac{F_{\mathrm{req}}}{\mu_s} = \frac{\qty{33.33}{\N} }{\num{0.3} } = \qty{111,11}{\N}
  .\end{equation*}
  To eliminate the reactions at $A$ we take the moment equilibrium about $A$ and get:
  \begin{align*}
    M_A &= 0 \\
    0 &= \qty{600}{\mm} \cdot P + \qty{50}{\mm} \cdot F_{\mathrm{req}} - \qty{200}{\mm} \cdot N \\
    P &= \frac{\qty{200}{\mm} \cdot N - \qty{50}{\mm} \cdot F_{\mathrm{req}}}{\qty{600}{\mm} } \\
    P &= \frac{\qty{200}{\mm} \cdot \qty{111.11}{\N} - \qty{50}{\mm} \cdot \qty{33.33}{\N} }{\qty{600}{\mm} }  \\
    P &= \qty{34,23}{\N}
  .\end{align*}
  Hence $P \geq \qty{34,23}{\N}$ and (a) $P = \qty{30}{\N}$ is therefore not enough whilst (b) $P = \qty{70}{\N}$ is enough.
\end{derivation}

\finalresult{
\begin{aligned}
  &(\mathrm{a}) & &\text{Brake does not hold}, \\
  &(\mathrm{b}) & &\text{Brake does hold}.
\end{aligned}
}
